\begin{trapbox}

	\begin{description}[style=nextline, leftmargin=2.8em, labelsep=0.5em, itemsep=0.35em]
		\item[\textbf{\textcolor{CTrap}{易错点一（未定式误判）}}]
			看见极限 $\frac{f(x)}{g(x)}$ 分子分母都趋近 $0$ 或 $\infty$ 就直接用洛必达，不检查是否真正满足条件（如 $g'(x)\neq0$、导数比极限存在等）。
		\item[\textbf{\textcolor{CTrap}{正确理解（条件审查）：}}]
			使用前必须确认：(1) $\lim f,\lim g$ 同时为 $0$ 或同时为 $\infty$；(2) 在去心邻域内 $f,g$ 可导，$g'(x)\neq0$ 且 $g(x)\neq0$；(3) $\lim f'/g'$ 存在（或无穷）。
		\item[\textbf{\textcolor{CTrap}{错因分析（粗心疏忽）：}}]
			对洛必达法则的适用条件记忆不完整，以为只要形式是 $\frac{0}{0}$ 或 $\frac{\infty}{\infty}$ 就可以直接套用，忽略了可导及极限存在的前提。
	\end{description}

	\medskip
	\begin{description}[style=nextline, leftmargin=2.8em, labelsep=0.5em, itemsep=0.35em]
		\item[\textbf{\textcolor{CTrap}{易错点二（必要当充分）}}]
			在恒成立问题中，用洛必达求出端点极限后就直接写成 $a\le \lim h(x)$ 或 $a\ge \lim h(x)$ 作为最终答案，未验证充分性。
		\item[\textbf{\textcolor{CTrap}{正确理解（最值思想）：}}]
			端点极限只是必要条件，必须进一步研究 $h(x)$ 的值域或最值，或证明在所得范围下不等式恒成立，才能确认最终范围。
		\item[\textbf{\textcolor{CTrap}{错因分析（思维定势）：}}]
			误认为极限值就是函数的最值，忽略了函数可能在区间内部取到更极端值；应牢记“必要性+充分性验证=完备”的解题模式。
	\end{description}

	\medskip
	\begin{description}[style=nextline, leftmargin=2.8em, labelsep=0.5em, itemsep=0.35em]
		\item[\textbf{\textcolor{CTrap}{易错点三（忽略极限存在）}}]
			直接使用洛必达，但 $\lim\frac{f'(x)}{g'(x)}$ 不存在（比如震荡）时，却错误地认为原极限也不存在。
		\item[\textbf{\textcolor{CTrap}{正确理解（失效情形）：}}]
			若 $\lim\frac{f'(x)}{g'(x)}$ 不存在（非无穷），不能判断原极限是否存在，此时洛必达法则失效，需改用其他方法。
		\item[\textbf{\textcolor{CTrap}{错因分析（逻辑倒置）：}}]
			误以为洛必达法则的逆命题成立，即若导数比极限不存在，则原极限必不存在；实际上 $\lim f'/g'$ 没有极限时原极限仍可能存在（例如 $f(x)=x^2\sin\frac1x,g(x)=x$ 在 $x\to0$）。
	\end{description}

\end{trapbox}
